% Source: https://artofproblemsolving.com/community/c3h1463151_a_stupid_question_i_had
% Stable identity: aops-topic-1463151-q01
Call the balls $a$, $b$, $c$. The student's count lists \emph{ordered} picks:
\[
(a,b),\ (b,a),\ (a,c),\ (c,a),\ (b,c),\ (c,b).
\]
But choosing $2$ balls only asks \emph{which} two balls we end up with, and the picks $(a,b)$ and $(b,a)$ give the same pair $\{a,b\}$. Every pair is listed exactly twice (once in each order), so the number of pairs is $\dfrac{6}{2}=\boxed{3}$, namely $\{a,b\}$, $\{a,c\}$, $\{b,c\}$.

The number $6=P(3,2)$ counts permutations: ordered selections, such as choosing a first prize and a second prize. In general $\dbinom{n}{k}=\dfrac{P(n,k)}{k!}$, since each set of $k$ objects can be put in order in $k!$ ways.
